\chapter{Limita\c{c}\~oes, riscos e evolu\c{c}\~ao}
\label{cap:limitacoes}

\section{Limita\c{c}\~oes t\'ecnicas}

\begin{table}[H]
\centering
\caption{Limita\c{c}\~oes por \'area e consequ\^encias.}
\label{tab:limitacoes}
\begin{tabularx}{\linewidth}{@{}p{3.3cm}X@{}}
\toprule
\textbf{\'Area} & \textbf{Limita\c{c}\~ao e consequ\^encia} \\
\midrule
Interrup\c{c}\~oes & IDT com 256 portas, PIC remapeado e timer IRQ0. IRQ1
do teclado continua mascarada; n\~ao h\'a APIC nem SMP, e a IRQ0 alimenta o
escalonador preemptivo. \\
Mem\'oria & Alocador de frames, p\'aginas de 4 KiB sob demanda e um heap
\emph{first-fit} de 1 MiB com \code{Vec}/\code{String}. Sem slab ou buddy
allocator, sem swap, sem reclaim e sem NUMA. \\
Processos & Tabela de tarefas, PID/PPID, troca de contexto, escalonador
preemptivo justo e ponderado com \code{nice}, pesos no estilo
\code{sched\_prio\_to\_weight} e \code{vruntime}, \code{sleep} bloqueante
(\code{schedule\_timeout}), filas de espera por evento
(\code{wait\_event}/\code{wake\_up}), sinais no estilo Unix com m\'ascaras de
pendentes/bloqueados e disposi\c{c}\~oes, tarefas paradas/continuadas, status
de sa\'ida, zumbis e \code{wait}/\code{SIGCHLD}. N\~ao h\'a
\code{fork}/\code{exec}, \'arvore rubro-negra completa do CFS/EEVDF nem
espa\c{c}o de endere\c{c}o por processo; os programas em ring 3 ainda n\~ao
s\~ao processos escalonados. Doze syscalls
(\code{exit/write/read/get\_tty/open/close/read\_fd/write\_fd/lseek/
getpid/getppid/kill}). \\
Armazenamento & VFS com dezesseis descritores sobre o ramfs em mem\'oria e um
\code{/proc} sint\'etico. Sem driver de disco, block layer, mounts,
permiss\~oes, page cache, FAT ou ext2; a ISO \'e s\'o o meio de boot. \\
Rede & Sem pilha TCP/IP, netdevice ou drivers de rede. \\
V\'ideo & Apenas VGA texto legado. UEFI sem CSM, GPUs e modos gr\'aficos
n\~ao s\~ao suportados. \\
Entrada & Teclado PS/2 por polling, sem USB-HID nativo, Caps Lock, AltGr,
setas ou layout ABNT2. \\
Serial & Apenas COM1, polling, sem IRQ, DMA ou controle de fluxo. \\
Concorr\^encia & Uma CPU; o escalonador \'e preemptivo por timer com
contabilidade justa e ponderada por \code{nice}/\code{vruntime},
\code{sleep} bloqueante e filas de espera por evento, mas a preemp\c{c}\~ao
fica suspensa enquanto um programa de ring 3 est\'a ativo. Os spinlocks usam
\code{irqsave}, ent\~ao a tabela de tarefas \'e segura contra a IRQ0. \\
Seguran\c{c}a & U/S separa usu\'ario e kernel, mas n\~ao h\'a ASLR, NX,
SMEP/SMAP, \emph{stack canary} nem espa\c{c}os de endere\c{c}o por processo.
Ponteiros do usu\'ario s\~ao validados por intervalo, mas n\~ao h\'a um
auxiliar no estilo \code{copy\_from\_user}. \\
Robustez & IDT evita \emph{triple fault} na maioria dos casos; n\~ao h\'a
recupera\c{c}\~ao de falhas, watchdog ou logs persistentes. \\
Portabilidade & C\'odigo espec\'ifico para x86\_64 BIOS; n\~ao compila para
ARM, RISC-V ou UEFI puro. \\
\bottomrule
\end{tabularx}
\end{table}

\section{Riscos}

\begin{table}[H]
\centering
\caption{Riscos e mitiga\c{c}\~oes atuais.}
\label{tab:riscos}
\begin{tabularx}{\linewidth}{@{}p{3.7cm}p{3.0cm}X@{}}
\toprule
\textbf{Risco} & \textbf{Probabilidade} & \textbf{Mitiga\c{c}\~ao} \\
\midrule
\emph{Triple fault} por exce\c{c}\~ao & M\'edia em c\'odigo novo & IDT e
\emph{stubs} j\'a instalados; testar novos tratadores em QEMU antes de
habilitar rotas cr\'iticas. \\
\emph{Deadlock} no console & Baixa hoje & N\~ao adquirir o lock duas vezes;
o timer n\~ao imprime, mas futuras IRQs devem usar \code{irqsave}. \\
\emph{Hang} em espera de hardware & M\'edia & Timeouts nas esperas do 8042;
a UART ainda pode travar se presente e defeituosa. \\
Teclado n\~ao funciona em hardware real & M\'edia & Testar com \emph{legacy
USB}/PS/2; usar serial como alternativa. \\
Tela preta em UEFI & Alta sem CSM & Documentar requisito de BIOS/CSM;
implementar GOP no futuro. \\
Perda de bytes na serial & M\'edia & Polling r\'apido; futuramente FIFO/IRQ. \\
Corrup\c{c}\~ao por bug de \'indice & Baixa & Revisar invariantes de
\code{row}/\code{column}; adicionar testes. \\
\bottomrule
\end{tabularx}
\end{table}

\section{Roadmap sugerido}

A ordem abaixo prioriza ganhos de robustez e de compreens\~ao antes de
adicionar complexidade.

\begin{table}[H]
\centering
\caption{Roadmap priorizado.}
\label{tab:roadmap}
\begin{tabularx}{\linewidth}{@{}clX@{}}
\toprule
\textbf{Prior.} & \textbf{Item} & \textbf{Objetivo} \\
\midrule
1 & Sinais em ring 3 e pilhas por tarefa & A entrega de sinais j\'a funciona
para tarefas de kernel; estend\^e-la a programas de ring 3 e processos
escalonados com pilha de kernel por tarefa. \\
2 & Processos de verdade & \code{fork}/\code{exec} (com \code{wait} j\'a
presente), tabelas de p\'agina por processo e tarefas de ring 3 escalonadas
pelo kernel. \\
3 & VFS e sistema de arquivos & Inodes, dentries, mounts e permiss\~oes;
driver de disco e initramfs. \\
4 & Robustez de IRQ e teclado & IRQ1 com buffer de scancodes e tratamento de
IRQs esp\'urias. \\
5 & Disciplina de linha TTY & Modo can\^onico, sinais, controle de jobs,
\code{Ctrl+C}/\code{Ctrl+D} e um \code{getty} por console. \\
6 & Framebuffer GOP/VBE & Suporte a UEFI e modos gr\'aficos; fonte bitmap. \\
7 & Endurecimento de mem\'oria & NX, SMEP/SMAP, valida\c{c}\~ao de ponteiros
de usu\'ario e ASLR. \\
8 & Driver de teclado completo & USB-HID, Caps Lock, AltGr, layout ABNT2,
teclas estendidas. \\
9 & Rede & NE2000/virtio-net, ARP/IP/ICMP/UDP, testes com QEMU. \\
10 & Testes automatizados & CI com QEMU headless, dumps VGA e asserts. \\
\bottomrule
\end{tabularx}
\end{table}

\section{Crit\'erios para considerar o projeto ``um kernel''}

Do ponto de vista did\'atico, o projeto j\'a \'e um kernel: roda em ring 0,
controla hardware e n\~ao depende de sistema operacional. Do ponto de vista
de um sistema operacional utiliz\'avel, ainda faltam, no m\'inimo:

\begin{itemize}
  \item isolamento de processos e espa\c{c}o de usu\'ario por processo;
  \item entrega de sinais em programas de ring 3, espa\c{c}os de endere\c{c}o
        por processo e pilhas de kernel por tarefa para ring 3
        (\code{nice}, pesos, \code{vruntime}, sleep, filas de espera por
        evento, sinais em contexto de kernel e \code{wait} j\'a existem);
  \item sistema de arquivos persistente e \emph{loader} de programas;
  \item drivers de entrada/sa\'ida al\'em de VGA/PS/2/serial;
  \item tratamento de erros e recupera\c{c}\~ao.
\end{itemize}

O projeto atual \'e melhor classificado como \textbf{kernel monol\'itico
m\'inimo de demonstra\c{c}\~ao}, focado em boot, mem\'oria,
interrup\c{c}\~oes, arquivos, tarefas e console.
